% ============================================================
% Chapter 6: Experimental Design
%   Section: Construction of the Custom Validation Set
%     Subsection: Limitations
% Included from chapters/06-experimental-design/03-custom-validation-set/custom-validation-set.tex
% ============================================================

\subsection{Limitations}
\label{subsec:custom-limitations}

\begin{itemize}
    \item \textbf{One annotator.} Self-agreement measures the
    consistency of the guide's application, not its validity: a
    systematic misreading applied twice agrees with itself.
    \item \textbf{Selection by the system.} Articles are drawn at
    random and claims are proposed by the system, which removes the
    annotator's tendency to pick the most checkable cases, but ties
    the set to the claims the system's selector favours. A claim the
    selector would never propose cannot enter the set, so the
    evaluation measures verification on the system's own claim
    choice, not on an independent sample of all check-worthy claims.
    The first six records were labelled before the daily draw existed
    and were chosen by hand; they are identified as such (no batch
    refers to them) and reported separately.
    \item \textbf{Temporal asymmetry.} Rule~3 restricts the
    annotator to evidence available at publication time, whereas the
    system searches today's web, without a date restriction. A system
    verdict informed by later evidence can therefore be scored as an
    error; such cases are identified from the evidence dates and
    reported separately rather than silently absorbed into the error
    rate.
    \item \textbf{Size.} As computed in
    Section~\ref{subsec:custom-size}, differences of a few points
    between models are within the uncertainty of a 150-claim set.
\end{itemize}
